\documentclass[10pt]{article}
\usepackage[T1]{fontenc}
\usepackage[a4paper,margin=19mm]{geometry}
\usepackage{amsmath,amssymb,booktabs,array}
\usepackage{enumitem}
\usepackage{xspace}
\usepackage{xcolor}
\usepackage[colorlinks=true,linkcolor=blue!45!black,urlcolor=blue!45!black]{hyperref}
\definecolor{accent}{HTML}{315B7D}
\setlength{\parindent}{0pt}
\setlength{\parskip}{4pt}
\setlist{nosep,leftmargin=*}
\newcommand{\goal}[1]{\medskip\noindent\textcolor{accent}{\textbf{Experimental goal.}} #1\medskip}
\newcommand{\code}[1]{\texttt{#1}}
\newcommand{\Wten}{\ifmmode\mathrm{W10}\else W10\fi}
\pagestyle{plain}

\title{TimeTensors Experiment 1: Constant Filtering}
\author{}
\date{}
\begin{document}\maketitle
\section{Theoretical overview and goal}
A constant look-back has zero empirical scale and carries little dynamic
information.  It can make normalized losses ill-conditioned and may correspond
to missing-value imputation rather than a forecasting regime.  Removing only
constant windows changes the sample distribution; dropping every affected user
also changes the population.

\goal{Separate numerical/optimization benefits of excluding constant training
windows from the easier evaluation obtained by excluding difficult test
windows or entire users.}

\section{Implementation details}
A look-back is flagged when it is constant or non-finite.  Window removal and
user removal are independently controlled for train and evaluation, yielding
the keep, train-only, evaluation-only, and both-split comparisons.  User removal
drops a user if any accessible look-back in the selected split is flagged.

Raw dataset preparation first reads the sibling \code{config.json}. Portable
top-level options, \code{timetensors}-scoped options, and run-specific additions
share one precedence rule; curated \code{drop\_users} lists are additive. These
source exclusions precede, and are distinct from, the dynamic constant-window
policies compared here.

The default study uses six datasets; $(168,24)$, $(504,168)$, and $(504,504)$;
DLinear; seeds 1--3; random sampling; batch size 256; learning rate $10^{-5}$;
10,000 optimizer steps; and validation every 1,000 steps. It compares keep,
drop-training-users, and drop-all-users first because prior experiments suggest
that window removal adds little after affected users are dropped. The full
profile retains all seven filtering policies. The launcher is
\code{01\_constants.slurm}.

Tables report a fixed test metric, seed mean $\pm$ sample standard deviation,
and explicit per-row $\times10^m$ scaling.  The number of retained windows and
users must accompany results: otherwise improvements caused by evaluating a
smaller population are indistinguishable from model improvements.

\section{Interpretation}
Train-only filtering is the cleanest test of optimization benefit.  Evaluation
filtering and user dropping answer different population questions and should
not be ranked directly against the unfiltered test without reporting coverage.
\end{document}
